\anexossection
\addcontentsline{toc}{section}{ANEXOS} 
\markboth{ANEXOS}
\begin{Glosario}
    \centering
    \caption{Glosario de términos}
    \label{tab:glosario-terminos}

    \footnotesize
    \setlength{\tabcolsep}{6pt}
    \renewcommand{\arraystretch}{1.25}

    \begin{tabular}{
        p{0.20\textwidth}
        p{0.72\textwidth}
    }
        \toprule
        \textbf{Término}
        & \textbf{Definición} \\
        \midrule

        CIE-10 (ICD-10)
        & Clasificación Internacional de Enfermedades, décima revisión.
          Estándar de la OMS para codificación de diagnósticos médicos. \\

        CPT
        & \textit{Current Procedural Terminology}. Sistema de codificación
          de procedimientos médicos de la AMA. \\

        EPS
        & Entidad Prestadora de Salud. Organización que administra planes
          de salud complementarios en el Perú. \\

        FHIR
        & \textit{Fast Healthcare Interoperability Resources}. Estándar de
          interoperabilidad de datos clínicos de HL7 International. \\

        \textit{Golden Set}
        & Conjunto curado de liquidaciones con clasificaciones esperadas,
          utilizado para regresión y calibración. \\

        IAFAS
        & Institución Administradora de Fondos de Aseguramiento en Salud.
          Entidad que administra fondos de seguro de salud. \\

        IPRESS
        & Institución Prestadora de Servicios de Salud. Clínica u hospital
          que brinda atención y presenta liquidaciones. \\

        LLM
        & \textit{Large Language Model}. Modelo de lenguaje de gran escala
          entrenado sobre corpus textuales masivos. \\

        MLR
        & \textit{Medical Loss Ratio}. Razón entre siniestros pagados y primas
          cobradas; indicador de eficiencia de la aseguradora. \\

        MVP
        & \textit{Minimum Viable Product}. Versión mínima funcional del producto
          que permite validar hipótesis de negocio. \\

        PII
        & \textit{Personally Identifiable Information}. Datos que permiten
          identificar a una persona natural. \\

        SUSALUD
        & Superintendencia Nacional de Salud. Ente regulador del sistema
          de salud en el Perú. \\

        TEDEF
        & Plataforma de intercambio electrónico de datos entre EPS e IPRESS
          para liquidaciones y acreditación. \\

        \bottomrule
    \end{tabular}
\end{table}


    